\documentclass[10pt,a4paper]{amsart}
\usepackage{amsmath,amssymb,amsthm,graphicx}
% Retained preamble for the Stacks Project chapter snapshot in this directory.
% The chapter's own preamble.tex is a copy of the shared Stacks Project preamble; the block
% below is sliced MECHANICALLY, line for line, from it. Nothing was added, removed or
% redefined; the wrapper class, the xr-hyper/hyperref wiring and the externaldocument list
% are omitted because the assembler supplies the document class.
%
% lines taken from preamble.tex: 3, 5, 9, 11, 12, 156, 157, 158, ... (51 lines total)
%
% package coverage restored from the chapter's own preamble (2026-10-13): xy, xy-2cell

\usepackage[all]{xy}
\xyoption{2cell}
\UseAllTwocells

\newenvironment{reference}{\comment}{\endcomment}
\newenvironment{slogan}{\comment}{\endcomment}
\newenvironment{history}{\comment}{\endcomment}
\theoremstyle{plain}
\newtheorem{theorem}[subsection]{Theorem}
\newtheorem{proposition}[subsection]{Proposition}
\newtheorem{lemma}[subsection]{Lemma}
\theoremstyle{definition}
\newtheorem{definition}[subsection]{Definition}
\newtheorem{example}[subsection]{Example}
\newtheorem{exercise}[subsection]{Exercise}
\newtheorem{situation}[subsection]{Situation}
\theoremstyle{remark}
\newtheorem{remark}[subsection]{Remark}
\newtheorem{remarks}[subsection]{Remarks}
\numberwithin{equation}{subsection}
\def\lim{\mathop{\mathrm{lim}}\nolimits}
\def\colim{\mathop{\mathrm{colim}}\nolimits}
\def\Spec{\mathop{\mathrm{Spec}}}
\def\Hom{\mathop{\mathrm{Hom}}\nolimits}
\def\Ext{\mathop{\mathrm{Ext}}\nolimits}
\def\SheafHom{\mathop{\mathcal{H}\!\mathit{om}}\nolimits}
\def\SheafExt{\mathop{\mathcal{E}\!\mathit{xt}}\nolimits}
\def\Sch{\mathit{Sch}}
\def\Mor{\mathop{\mathrm{Mor}}\nolimits}
\def\Ob{\mathop{\mathrm{Ob}}\nolimits}
\def\Sh{\mathop{\mathit{Sh}}\nolimits}
\def\NL{\mathop{N\!L}\nolimits}
\def\CH{\mathop{\mathrm{CH}}\nolimits}
\def\proetale{{pro\text{-}\acute{e}tale}}
\def\etale{{\acute{e}tale}}
\def\QCoh{\mathit{QCoh}}
\def\Ker{\mathop{\mathrm{Ker}}}
\def\Im{\mathop{\mathrm{Im}}}
\def\Coker{\mathop{\mathrm{Coker}}}
\def\Coim{\mathop{\mathrm{Coim}}}
\DeclareMathSymbol{\boxtimes}{\mathbin}{AMSa}{"02}
\def\QCohstack{\mathcal{QC}\!\mathit{oh}}
\def\Cohstack{\mathcal{C}\!\mathit{oh}}
\def\Spacesstack{\mathcal{S}\!\mathit{paces}}
\def\Quotfunctor{\mathrm{Quot}}
\def\Hilbfunctor{\mathrm{Hilb}}
\def\Curvesstack{\mathcal{C}\!\mathit{urves}}
\def\Polarizedstack{\mathcal{P}\!\mathit{olarized}}
\def\Complexesstack{\mathcal{C}\!\mathit{omplexes}}
\def\Pic{\mathop{\mathrm{Pic}}\nolimits}
\def\Picardstack{\mathcal{P}\!\mathit{ic}}
\def\Picardfunctor{\mathrm{Pic}}
\def\Deformationcategory{\mathcal{D}\!\mathit{ef}}
\usepackage{cleveref}
\crefname{originallemma}{original lemma}{original lemmas}
\Crefname{originallemma}{Original lemma}{Original lemmas}
\makeatletter
\newlabel{equation-extension-ringed-topoi}{{08MB}{source}{}{sourcepointer-1}{}}
\newlabel{equation-extension-ringed-topoi@cref}{{[originallemma][0][]08MB}{[1][0][]source}{}{}{}}
\makeatother
\crefname{originallemma}{original lemma}{original lemmas}
\Crefname{originallemma}{Original lemma}{Original lemmas}
\makeatletter
\newlabel{lemma-inf-ext-ringed-topoi}{{08ME}{source}{}{sourcepointer-2}{}}
\newlabel{lemma-inf-ext-ringed-topoi@cref}{{[originallemma][0][]08ME}{[1][0][]source}{}{}{}}
\makeatother
\crefname{originallemma}{original lemma}{original lemmas}
\Crefname{originallemma}{Original lemma}{Original lemmas}
\makeatletter
\newlabel{sites-modules-lemma-module-quotient-flat}{{03EW}{source}{}{sourcepointer-3}{}}
\newlabel{sites-modules-lemma-module-quotient-flat@cref}{{[originallemma][0][]03EW}{[1][0][]source}{}{}{}}
\makeatother
\title{Stacks, Deformation Theory: lemma-inf-obs-ext-ringed-topoi}
\author{The Stacks Project Authors}
\date{stacks-project:defos.tex:local-snapshot-20261012; GNU FDL 1.2 or later}
\begin{document}
\maketitle
\noindent\textit{Editorial note.} Statement span is lines 2946-2961 of sources/stacks-defos/pinned-original-source.tex; proof span is lines 2963-3077, the complete author proof (113 body lines).\par\medskip
\noindent\textit{Editorial note.} Navigation stubs: 3 cited-but-not-included labels are supplied as non-load-bearing pointers with their pinned tags where the snapshot tag index has them.\par\medskip
\noindent\textit{Editorial note.} extra\_packages is empty on purpose: the retained preamble already loads xy with options and a second declaration clashes in the assembler wrapper.\par\medskip
\noindent\textit{Editorial note.} Text is retained verbatim, including the author's own typos and placeholders.\par\medskip
\noindent\textit{Editorial note.} Queue index 3 of queue\_defos.json (pb=3982, stmt 2946-2961, proof 2963-3077).\par\medskip
\noindent\textit{Allowed standard prerequisites (editorial).}\begin{itemize}
\item Explicit prerequisite boundary. The cited inputs are prerequisites; the argument in the retained span is not replaced by any cited result.
\end{itemize}\medskip
\ifdefined\hypertarget\hypertarget{sourcepointer-1}{}\fi
\noindent\textit{Source reference.} Original lemma 08MB: sources/stacks-defos/pinned-original-source.tex, source lines 2770--2772. Cited by the retained statement or proof but not part of the collected spans; supplied only as a navigation pointer. The argument does not depend on its content.\par\medskip
\ifdefined\hypertarget\hypertarget{sourcepointer-2}{}\fi
\noindent\textit{Source reference.} Original lemma 08ME: sources/stacks-defos/pinned-original-source.tex, source lines 2912--2922. Cited by the retained statement or proof but not part of the collected spans; supplied only as a navigation pointer. The argument does not depend on its content.\par\medskip
\ifdefined\hypertarget\hypertarget{sourcepointer-3}{}\fi
\noindent\textit{Source reference.} Original lemma 03EW: sources/stacks-defos/pinned-original-source.tex, source lines 3077--3077. Cited by the retained statement or proof but not part of the collected spans; supplied only as a navigation pointer. The argument does not depend on its content.\par\medskip
\begin{lemma}
\label{lemma-inf-obs-ext-ringed-topoi}
Let $i : (\Sh(\mathcal{C}), \mathcal{O}) \to (\Sh(\mathcal{D}), \mathcal{O}')$
be a first order thickening of ringed topoi. Assume given
$\mathcal{O}$-modules $\mathcal{F}$, $\mathcal{K}$
and an $\mathcal{O}$-linear map
$c : \mathcal{I} \otimes_\mathcal{O} \mathcal{F} \to \mathcal{K}$.
Then there exists an element
$$
o(\mathcal{F}, \mathcal{K}, c) \in
\Ext^2_\mathcal{O}(\mathcal{F}, \mathcal{K})
$$
whose vanishing is a necessary and sufficient condition for the existence
of a sequence (\ref{equation-extension-ringed-topoi})
with $c_{\mathcal{F}'} = c$.
\end{lemma}

\begin{proof}
We first show that if $\mathcal{K}$ is an injective $\mathcal{O}$-module,
then there does exist a sequence (\ref{equation-extension-ringed-topoi}) with
$c_{\mathcal{F}'} = c$. To do this, choose a flat
$\mathcal{O}'$-module $\mathcal{H}'$ and a surjection
$\mathcal{H}' \to \mathcal{F}$
(Modules on Sites, Lemma \ref{sites-modules-lemma-module-quotient-flat}).
Let $\mathcal{J} \subset \mathcal{H}'$ be the kernel. Since $\mathcal{H}'$
is flat we have
$$
\mathcal{I} \otimes_{\mathcal{O}'} \mathcal{H}' =
\mathcal{I}\mathcal{H}'
\subset \mathcal{J} \subset \mathcal{H}'
$$
Observe that the map
$$
\mathcal{I}\mathcal{H}' =
\mathcal{I} \otimes_{\mathcal{O}'} \mathcal{H}'
\longrightarrow
\mathcal{I} \otimes_{\mathcal{O}'} \mathcal{F} =
\mathcal{I} \otimes_\mathcal{O} \mathcal{F}
$$
annihilates $\mathcal{I}\mathcal{J}$. Namely, if $f$ is a local section
of $\mathcal{I}$ and $s$ is a local section of $\mathcal{H}$, then
$fs$ is mapped to $f \otimes \overline{s}$ where $\overline{s}$ is
the image of $s$ in $\mathcal{F}$. Thus we obtain
$$
\xymatrix{
\mathcal{I}\mathcal{H}'/\mathcal{I}\mathcal{J}
\ar@{^{(}->}[r] \ar[d] &
\mathcal{J}/\mathcal{I}\mathcal{J} \ar@{..>}[d]_\gamma \\
\mathcal{I} \otimes_\mathcal{O} \mathcal{F} \ar[r]^-c &
\mathcal{K}
}
$$
a diagram of $\mathcal{O}$-modules. If $\mathcal{K}$ is injective
as an $\mathcal{O}$-module, then we obtain the dotted arrow.
Denote $\gamma' : \mathcal{J} \to \mathcal{K}$ the composition
of $\gamma$ with $\mathcal{J} \to \mathcal{J}/\mathcal{I}\mathcal{J}$.
A local calculation shows the pushout
$$
\xymatrix{
0 \ar[r] &
\mathcal{J} \ar[r] \ar[d]_{\gamma'} &
\mathcal{H}' \ar[r] \ar[d] &
\mathcal{F} \ar[r] \ar@{=}[d] &
0 \\
0 \ar[r] &
\mathcal{K} \ar[r] &
\mathcal{F}' \ar[r] &
\mathcal{F} \ar[r] &
0
}
$$
is a solution to the problem posed by the lemma.

\medskip\noindent
General case. Choose an embedding $\mathcal{K} \subset \mathcal{K}'$
with $\mathcal{K}'$ an injective $\mathcal{O}$-module. Let $\mathcal{Q}$
be the quotient, so that we have an exact sequence
$$
0 \to \mathcal{K} \to \mathcal{K}' \to \mathcal{Q} \to 0
$$
Denote
$c' : \mathcal{I} \otimes_\mathcal{O} \mathcal{F} \to \mathcal{K}'$
be the composition. By the paragraph above there exists a sequence
$$
0 \to \mathcal{K}' \to \mathcal{E}' \to \mathcal{F} \to 0
$$
as in (\ref{equation-extension-ringed-topoi}) with $c_{\mathcal{E}'} = c'$.
Note that $c'$ composed with the map $\mathcal{K}' \to \mathcal{Q}$
is zero, hence the pushout of $\mathcal{E}'$ by
$\mathcal{K}' \to \mathcal{Q}$ is an extension
$$
0 \to \mathcal{Q} \to \mathcal{D}' \to \mathcal{F} \to 0
$$
as in (\ref{equation-extension-ringed-topoi}) with $c_{\mathcal{D}'} = 0$.
This means exactly that $\mathcal{D}'$ is annihilated by
$\mathcal{I}$, in other words, the $\mathcal{D}'$ is an extension
of $\mathcal{O}$-modules, i.e., defines an element
$$
o(\mathcal{F}, \mathcal{K}, c) \in
\Ext^1_\mathcal{O}(\mathcal{F}, \mathcal{Q}) =
\Ext^2_\mathcal{O}(\mathcal{F}, \mathcal{K})
$$
(the equality holds by the long exact cohomology sequence associated
to the exact sequence above and the vanishing of higher ext groups
into the injective module $\mathcal{K}'$). If
$o(\mathcal{F}, \mathcal{K}, c) = 0$, then we can choose a splitting
$s : \mathcal{F} \to \mathcal{D}'$ and we can set
$$
\mathcal{F}' = \Ker(\mathcal{E}' \to \mathcal{D}'/s(\mathcal{F}))
$$
so that we obtain the following diagram
$$
\xymatrix{
0 \ar[r] &
\mathcal{K} \ar[r] \ar[d] &
\mathcal{F}' \ar[r] \ar[d] &
\mathcal{F} \ar[r] \ar@{=}[d] &
0 \\
0 \ar[r] &
\mathcal{K}' \ar[r] &
\mathcal{E}' \ar[r] &
\mathcal{F} \ar[r] & 0
}
$$
with exact rows which shows that $c_{\mathcal{F}'} = c$. Conversely, if
$\mathcal{F}'$ exists, then the pushout of $\mathcal{F}'$ by the map
$\mathcal{K} \to \mathcal{K}'$ is isomorphic to $\mathcal{E}'$ by
Lemma \ref{lemma-inf-ext-ringed-topoi} and the vanishing of higher ext groups
into the injective module $\mathcal{K}'$. This gives a diagram
as above, which implies that $\mathcal{D}'$ is split as an extension, i.e.,
the class $o(\mathcal{F}, \mathcal{K}, c)$ is zero.
\end{proof}

\end{document}
